% Part: first-order-logic
% Chapter: models-theories
% Section: expressing-props-of-structures

\documentclass[../../../include/open-logic-section]{subfiles}

\begin{document}

\olfileid{fol}{mat}{exs}

\olsection{د \printtoken{P}{structure} د خاصيتونو بيانول}

\begin{explain}
د تابعو او اړيکو شرطونه، يا په لا عمومي ډول دا خبره بيانول چې په
يوه جوړښت کښې تابعې او اړيکې دغه شرطونه پوره کوي، ډېر ځله ګټور او
مهم وي۔ د بېلګې په توګه، غواړو د يوې ژبې هغه !!{structure}s چې
!!{predicate}s پکښې زموږ مطلوبه ``معنٰی'' لري، له هغو جوړښتونو څخه
بېل کړو چې داسې نۀ دي۔ البته بشپړه ازادي لرو چې خپل ``مطلوب''
!!{structure}s وټاکو؛ مثلاً ويلے شو چې د !!{predicate}~$\le$ تفسير بايد يو
ترتيب وي، يا يوازې د~$\Lang L$ له هغو تفسيرونو سره دلچسپي لرو چې
ساحه ئې له سټونو جوړه وي او~$\Obj \in$ د ``د !!a{element} په توګه
شتون'' په اړيکه تفسير شي۔ خو ايا دا د ژبې په !!{sentence}s کولے شو؟
په بل عبارت، پر !!a{structure}~$\Struct M$ کوم شرطونه د~$\Struct M$
د ژبې په يوه !!a{sentence} (يا ښايي د !!{sentence}s په يوه سټ)
بيانولے شو؟ ځينې شرطونه نۀ شو بيانولے۔ مثلاً د~$\Lang L_A$ داسې
هېڅ جمله نشته چې يوازې په هغه !!{structure}~$\Struct M$ کښې رښتونې
وي چې~$\Domain M = \Nat$ وي۔ ``ساحه يوازې طبيعي عددونه لري'' نۀ شو
بيانولے۔ خو ښايي د !!{structure}s ځينې ``جوړښتي خاصيتونه'' بيانولے
شو۔ د !!{structure}s کوم خاصيتونه په !!{sentence}s بيانولے شو؟ يا
په بله وينا، د !!{structure}s کومې ټولګې د هغو جوړښتونو په توګه
بيانولے شو چې يوه !!a{sentence} (يا د !!{sentence}s يو سټ) رښتونې
کوي؟
\end{explain}

\begin{defn}[د يوه سټ مدل]
فرض کړئ~$\Gamma$ په ژبه~$\Lang L$ کښې د !!{sentence}s يو سټ دے۔
وايو چې !!a{structure}~$\Struct M$ د~$\Gamma$ \emph{مدل دے}، که د
هر~$!A \in \Gamma$ دپاره $\Sat{M}{!A}$ وي۔
\end{defn}

\begin{ex}
جمله~$\lforall[x][x \le x]$ په~$\Struct M$ کښې هله او يوازې هله
رښتونې ده چې~$\Assign{\le}{M}$ انعکاسي اړيکه وي۔ جمله
$\lforall[x][\lforall[y][((x \le y \land y \le x) \lif x = y)]]$
په~$\Struct M$ کښې هله او يوازې هله رښتونې ده چې
$\Assign{\le}{M}$ ضد متناظره وي۔ جمله
$\lforall[x][\lforall[y][\lforall[z][((x \le y \land y \le z)
      \lif x \le z)]]]$ په~$\Struct M$ کښې هله او يوازې هله رښتونې
ده چې~$\Assign{\le}{M}$ انتقالي وي۔ نو د لاندې سټ مدلونه
\begin{align*}
\{\quad &\lforall[x][x \le x], \\
   & \lforall[x][\lforall[y][((x \le y \land y \le
    x) \lif x = y)]], \\
   &\lforall[x][\lforall[y][\lforall[z][((x \le y
      \land y \le z) \lif x \le z)]]] \quad \}
\end{align*}
کټ مټ هغه جوړښتونه دي چې~$\Assign{\le}{M}$ پکښې انعکاسي، ضد
متناظره او انتقالي، يعنې جزوي ترتيب وي۔ له همدې امله، دا د
\emph{جزوي ترتيبونو د لومړۍ درجې تيورۍ} د بديهي اصولو په توګه
اخستلے شو۔
\end{ex}

\end{document}
